\documentclass[11pt,twocolumn]{article}

% Geometry: tight 2-column with reasonable margins.
\usepackage[margin=0.75in]{geometry}
\usepackage{graphicx}
\usepackage{amsmath,amssymb}
\usepackage{booktabs}
\usepackage{xcolor}
\usepackage{url}
\usepackage[colorlinks=true,linkcolor=blue!50!black,citecolor=blue!50!black,urlcolor=blue!50!black]{hyperref}
\usepackage[numbers,sort&compress]{natbib}
\usepackage{authblk}
\usepackage{microtype}
\usepackage[section]{placeins}
\usepackage[skip=4pt]{caption}
\usepackage{flushend}
\graphicspath{{figures/}{../figures/}}

% Tighter float spacing - eliminate dead whitespace around figures and tables.
\setlength{\intextsep}{5pt plus 1pt minus 1pt}
\setlength{\textfloatsep}{6pt plus 1pt minus 1pt}
\setlength{\floatsep}{6pt plus 1pt minus 1pt}
\setlength{\dbltextfloatsep}{6pt plus 1pt minus 1pt}
\setlength{\dblfloatsep}{6pt plus 1pt minus 1pt}
\setlength{\abovecaptionskip}{3pt}
\setlength{\belowcaptionskip}{2pt}

% Boundary overflow protections - long identifiers and compact math.
\emergencystretch=2em
\tolerance=2000
\hyphenpenalty=50
\sloppy

% Slightly looser table rows for readability.
\renewcommand{\arraystretch}{1.10}

\title{Reuse versus Adaptation: Decomposing the ImageNet-to-CIFAR-10 Transfer Win for ResNet-18}
\author[1]{Vikash Chandra Mishra\thanks{\texttt{crimsonsyrus000@gmail.com}}}
\affil[1]{Independent}
\date{\today}

\begin{document}
\maketitle


\begin{abstract}
Transfer learning from ImageNet has become a default first move for image classification on small natural-image datasets, but the accuracy gain is often reported as a single number rather than decomposed into the value of \emph{reusing} pretrained features and the additional value of \emph{adapting} those features to the downstream task. We quantify both contributions on CIFAR-10 by training the same ResNet-18 backbone under three carefully controlled conditions: random initialisation with all weights trainable; ImageNet-pretrained weights with the backbone frozen and a linear probe on a new 10-class head; and ImageNet-pretrained weights with full end-to-end fine-tuning. All three conditions share an identical 5-epoch training schedule, batch size, and $224{\times}224$ input pipeline on the official CIFAR-10 train and test splits, so any difference in test accuracy is attributable to initialisation and parameter freedom alone. Across three random seeds, we observe a strict ordering on every individual seed: the from-scratch baseline reaches \textbf{0.7830} $\pm$ 0.0244 test top-1 accuracy, the linear probe over frozen ImageNet features reaches \textbf{0.8674} $\pm$ 0.0016, and full fine-tuning reaches \textbf{0.9119} $\pm$ 0.0025. The linear-probe gain over scratch isolates the value of frozen transfer at $+8.44$ percentage points; the additional gain from fine-tuning isolates the value of adaptation at $+4.45$ percentage points. We also find that the pretrained conditions are an order of magnitude more stable across seeds than the from-scratch run, consistent with pretrained initialisations occupying a flatter region of the loss landscape under tight training budgets. Code, logs, and per-seed numbers accompanying these claims are released as a public repository (URL on the project site).
\end{abstract}

\section{Introduction}

Image classification on small natural-image datasets is one of the
oldest and best-instrumented benchmarks in machine learning, yet the
practical advice given to practitioners working on a new dataset has
converged on a single recipe: take a backbone that has been trained on
ImageNet \cite{deng2009imagenet}, replace its final classification
layer with one sized to the new task, and fine-tune. The recipe works,
and works robustly, across an enormous range of downstream domains
\cite{kornblith2019better,yosinski2014transferable}. What is much less
often quantified, especially at a budget that a student or a
single-GPU laptop can reproduce, is \emph{why} it works. The accuracy
gain over training the same architecture from random initialisation
is typically reported as one number; that number conflates two
mechanisms that practitioners benefit from for very different reasons.

The first mechanism is the value of \emph{reusing} pretrained
features without further adaptation. If the frozen feature space that
ImageNet training produces is already approximately linearly
separable for a downstream class structure, a linear probe over those
frozen features captures essentially all of the easy signal at a
trainable-parameter cost three to four orders of magnitude smaller
than the full model. The second mechanism is the value of
\emph{adapting} those features to the downstream task: fine-tuning
moves the backbone away from its ImageNet optimum toward a downstream
optimum, trading transferability for specialisation. The relative
size of the two contributions matters in deployment, where freezing
the backbone gives a much smaller training footprint and a much
larger memory budget for batch size or input resolution; in
continual-learning scenarios, where backbone drift is a hazard; and
pedagogically, where students often see fine-tuning presented as a
single black-box operation rather than a composition of two
distinguishable effects.

\textbf{Contributions.} We isolate and quantify the two contributions
on CIFAR-10 \cite{krizhevsky2009learning}, the canonical small-image
benchmark, with the simplest experimental design that admits a clean
decomposition.
\begin{itemize}
  \item We train the same ResNet-18 \cite{he2016deep} backbone under
    three conditions that differ only in initialisation and which
    parameters are trainable: from scratch (random init, all weights
    trainable); a linear probe over a frozen ImageNet-pretrained
    backbone; and full end-to-end fine-tuning from the same
    ImageNet checkpoint. Every other ingredient (training schedule,
    batch size, input pipeline, optimiser per condition, three random
    seeds) is held fixed.
  \item Across three random seeds, the ordering Scratch $<$ Linear
    Probe $<$ Fine-tune holds on every individual seed and in the
    mean. Linear probing buys
    $+8.44$ percentage points of test
    accuracy over the from-scratch baseline; full fine-tuning buys a
    further $+4.45$ percentage points on top
    of the probe.
  \item The pretrained conditions exhibit per-seed standard
    deviations of $0.0016$ (linear probe) and $0.0025$
    (fine-tune), an order of magnitude tighter than the
    $0.0244$ of the from-scratch run, indicating
    that ImageNet initialisation does not just raise mean accuracy
    but also stabilises the optimisation trajectory under a tight
    training budget.
  \item We release per-seed metrics, the experiment script, and the
    LaTeX source of this paper as a self-contained, single-GPU
    reproduction package suitable for student-led
    transfer-learning study groups.
\end{itemize}

\textbf{Roadmap.} Section~\ref{sec:related} situates the
contributions in the broader transfer-learning literature.
Section~\ref{sec:method} formalises the three conditions and the
shared input pipeline. Section~\ref{sec:experiments} reports the
training protocol and dataset statistics. Section~\ref{sec:results}
gives the per-condition test-accuracy numbers, the per-epoch
trajectories, and the limitations that follow from a 5-epoch budget.
Section~\ref{sec:conclusion} discusses what these numbers do and do
not say about transfer learning more broadly.


\section{Related Work}
\label{sec:related}

\paragraph{ImageNet pretraining as a default initialiser.}
A long line of empirical work has shown that supervised pretraining
on ImageNet \cite{deng2009imagenet} produces backbones that transfer
robustly to a wide range of downstream image-classification tasks.
\citet{yosinski2014transferable} characterised feature transferability
across layers and showed that lower layers transfer more readily than
higher ones; \citet{kornblith2019better} found that better ImageNet
accuracy generally predicts better transfer, though with diminishing
returns. \citet{he2019rethinking} pushed back on the notion that
ImageNet pretraining is uniformly necessary, demonstrating that
sufficiently long training from scratch can match pretraining on some
detection benchmarks. The picture that emerges from these studies is
that ImageNet pretraining is a strong default at modest training
budgets but loses its edge when compute is plentiful, an observation
our 5-epoch results are designed to illustrate at the
single-laptop-GPU end of the budget spectrum.

\paragraph{Linear probing as a feature-quality probe.}
Linear probes over frozen pretrained features are a standard
diagnostic for whether a representation has captured the task's class
structure without requiring further adaptation. The technique
predates modern self-supervised learning but rose to prominence in
the self-supervised era, where it serves as a comparator across
different pretraining objectives \cite{chen2020simple,he2020momentum,
caron2020swav}. \citet{chen2020simple} popularised the practice of
reporting both the linear-probe accuracy and the fine-tuned accuracy
of the same backbone as separate numbers, since the two measure
different properties of the pretrained representation. Our
contribution sits in this tradition; we apply the same
probe-versus-fine-tune comparison to the supervised
ImageNet-pretrained backbone that practitioners reach for by default.

\paragraph{CIFAR-10 as a transfer-learning testbed.}
CIFAR-10 \cite{krizhevsky2009learning} has been used as a
transfer-learning destination since the earliest days of deep image
classification, both as a small target where overfitting is
particularly visible and as a sanity-check benchmark whose
state-of-the-art is well-characterised. The standard practice of
upsampling its native $32{\times}32$ inputs to the $224{\times}224$
resolution that ImageNet-trained backbones expect is documented in
\citet{kornblith2019better} and is a confound any transfer study on
CIFAR-10 must acknowledge: the upsampled CIFAR-10 image is not the
same statistical object as the native one, and absolute accuracies on
the two should not be compared directly.

\paragraph{Distinguishing the present work.}
Most existing transfer-learning studies on CIFAR-10 either compare
many backbones at one fine-tuning recipe \cite{kornblith2019better}
or compare many pretraining objectives at one downstream task
\cite{ericsson2021how}. We deliberately fix everything except
initialisation and parameter freedom, and we run the simplest
possible three-way decomposition (Scratch, Linear Probe, Fine-tune)
on a single backbone. The aim is not to set a new state of the art,
which would require longer training, hyperparameter search, and
modern augmentation \cite{cubuk2020randaugment}; it is to give a
clean, reproducible, single-GPU answer to the question of how much of
the transfer-learning win is feature reuse and how much is feature
adaptation.


\section{Method}
\label{sec:method}

We describe the three transfer-learning conditions, the shared
ResNet-18 backbone, and the shared CIFAR-10 input pipeline that the
conditions hold fixed. The motivation for fixing every variable
except initialisation and parameter freedom is that any difference in
test accuracy across conditions can then be attributed cleanly to one
of those two factors.

\subsection{Three transfer-learning conditions}

Let $f_\theta : \mathcal{X} \to \mathbb{R}^{512}$ denote the
ResNet-18 backbone parameterised by $\theta$, and let $g_\phi :
\mathbb{R}^{512} \to \mathbb{R}^{10}$ denote a randomly initialised
linear classifier with parameters $\phi$. We evaluate three training
conditions on the joint model $g_\phi \circ f_\theta$:
%
\begin{align*}
  &\text{Scratch:}      & \theta &\sim p_\text{rand},  & \nabla\theta \neq 0, \quad \nabla\phi \neq 0 \\
  &\text{Linear probe:} & \theta &= \theta_\text{IN},  & \nabla\theta = 0,    \quad \nabla\phi \neq 0 \\
  &\text{Fine-tune:}    & \theta &= \theta_\text{IN},  & \nabla\theta \neq 0, \quad \nabla\phi \neq 0
\end{align*}
%
Here $\theta_\text{IN}$ denotes the publicly released
\texttt{ResNet18\_Weights.IMAGENET1K\_V1} ImageNet-1K checkpoint
shipped with \texttt{torchvision} \cite{paszke2019pytorch}, and
$p_\text{rand}$ denotes its default random-init scheme. The
classifier head $\phi$ is randomly initialised in every condition,
including the Linear probe, so that the only thing the probe is
asked to learn is a 10-class projection of the frozen backbone's
final-layer features. Figure~\ref{fig:method-overview} renders the
three conditions side by side.

\begin{figure}[t]
  \centering
  \includegraphics[width=\linewidth]{figures/fig-method-overview.png}
  \caption{The three transfer-learning conditions evaluated in this
    paper. All three share the same ResNet-18 architecture and
    CIFAR-10 input pipeline; they differ only in the initialisation
    of the backbone and which parameters carry gradients.}
  \label{fig:method-overview}
\end{figure}

\subsection{Backbone}

The backbone $f_\theta$ is the standard 18-layer Residual Network
\cite{he2016deep} as implemented in \texttt{torchvision}: a
$7{\times}7$ stride-2 stem (Conv1, BatchNorm, ReLU, $3{\times}3$
max-pool), four residual stages each consisting of two BasicBlock
units, global average pooling, and a final 512-to-1000 linear layer.
We replace the final 512-to-1000 layer with a 512-to-10 linear layer
(the head $g_\phi$). The Linear-probe condition freezes
\emph{everything} below the head, including the BatchNorm running
statistics: every \texttt{BatchNorm2d} module is held in
\texttt{eval} mode through training. Without that pin, the running
means and variances would drift even when the convolution weights are
frozen, and the probe would no longer be measuring "frozen ImageNet
features" in any meaningful sense.

\subsection{Input pipeline}

CIFAR-10 \cite{krizhevsky2009learning} ships as $32{\times}32$ RGB
images, but ImageNet-pretrained backbones expect a much larger
receptive field. We therefore upsample every image to
$224{\times}224$ via bicubic interpolation before feeding it to the
backbone, and normalise with the ImageNet per-channel mean and
standard deviation that the pretrained checkpoint expects. The full
pipeline is shown in Figure~\ref{fig:data-pipeline}: training
augmentation is a single random horizontal flip on top of the resize,
and evaluation uses the resize and normalise stages alone. We do not
apply any heavier augmentation (CutMix, AutoAugment, MixUp) so that
the comparison across conditions is not confounded by augmentation
strength.

\begin{figure}[t]
  \centering
  \includegraphics[width=\linewidth]{figures/fig-data-pipeline.png}
  \caption{Shared CIFAR-10 input pipeline. The same upsample,
    normalise, and (training-only) flip steps are applied in every
    condition; only the downstream model differs.}
  \label{fig:data-pipeline}
\end{figure}

\subsection{Optimisation}

For the from-scratch and fine-tuning conditions we use Adam
\cite{kingma2015adam} with a fixed learning rate of $10^{-3}$. For
the linear probe we use SGD with momentum $0.9$ and a fixed learning
rate of $10^{-2}$, which is a standard choice for linear-probe
heads when the backbone supplies high-quality features
\cite{he2020momentum}. All three conditions train for five epochs
with batch size 128 and the standard cross-entropy loss; no
hyperparameter search is performed and no learning-rate schedule is
applied. Each seed in $\{0, 1, 2\}$ is run sequentially through all
three conditions, and we record per-epoch train and test metrics for
every (condition, seed) pair.


\section{Experiments}
\label{sec:experiments}

\subsection{Dataset}

We evaluate on CIFAR-10 \cite{krizhevsky2009learning}, a 60{,}000-image
ten-class dataset of $32{\times}32$ RGB photographs, split into the
official 50{,}000-image training set and 10{,}000-image test set. The
dataset is loaded from the \texttt{uoft-cs/cifar10} HuggingFace
release \cite{lhoest2021datasets}, which serves the standard image
bytes through a stable CDN; we observed intermittent 503 errors from
the canonical \texttt{www.cs.toronto.edu} mirror during this study,
and use the HuggingFace mirror as the single source. No data
augmentation other than a random horizontal flip during training is
applied (Section~\ref{sec:method}).

\subsection{Training protocol}

Each random seed in $\{0, 1, 2\}$ runs the three conditions
(Scratch, Linear Probe, Fine-tune) sequentially using the same
\texttt{torchvision} ResNet-18 architecture and the same
\texttt{ResNet18\_Weights.IMAGENET1K\_V1} pretrained checkpoint
(loaded only by the Linear Probe and Fine-tune conditions). All
conditions train for five epochs with batch size 128 and the standard
multi-class cross-entropy loss. The Scratch and Fine-tune conditions
optimise with Adam \cite{kingma2015adam} at learning rate $10^{-3}$;
the Linear Probe optimises only its 5{,}130 head parameters with SGD
at learning rate $10^{-2}$ and momentum $0.9$. We pin every
\texttt{BatchNorm2d} module to evaluation mode in the Linear Probe so
its running statistics do not drift during the forward-only passes.
Training is run on a single NVIDIA RTX A4000 (16\,GB) GPU and peaks
at 3.13\,GB of memory, well below the
card's capacity. The mean total wall-clock per seed across the three
conditions is 478\,s; the full three-seed
sweep finishes in well under thirty minutes of wall-clock time.

\subsection{Baselines and conditions}

The baseline structure of this study is unusual: rather than
comparing our method against external prior work, we compare three
in-run \emph{conditions} that differ only in initialisation and
parameter freedom, holding everything else fixed.
%
\begin{itemize}
  \item \textbf{Scratch.} ResNet-18 with random initialisation and
    all parameters trainable.
  \item \textbf{Linear Probe.} ResNet-18 backbone loaded from
    \texttt{ResNet18\_Weights.IMAGENET1K\_V1} and frozen
    (\texttt{requires\_grad}\,$=$\,\texttt{False} on every backbone
    parameter, BatchNorm pinned to \texttt{eval}); only the new
    512-to-10 classifier head is trained.
  \item \textbf{Fine-tune.} ResNet-18 backbone loaded from
    \texttt{ResNet18\_Weights.IMAGENET1K\_V1} and the new classifier
    head, all parameters trainable.
\end{itemize}
%
We deliberately do not include external CIFAR-10 baselines from the
literature in our headline numbers: the standard literature numbers
for ResNet-18 on CIFAR-10 use longer training schedules (often
100--200 epochs), heavier augmentation, and step or cosine
learning-rate schedules, none of which we apply. Reporting them
alongside our 5-epoch numbers would invite misreadings of the
contribution. We discuss the implications of this choice in
Section~\ref{sec:results}.


\subsection{Reproducibility}

To make the reported numbers reproducible at the byte level, we list
the exact configuration that produced them. Random seeds are
$\{0, 1, 2\}$ applied via \texttt{torch.manual\_seed}, \texttt{numpy.random.seed},
and Python's \texttt{random.seed} at the start of each (condition, seed)
pair. The pretrained checkpoint is the public
\texttt{ResNet18\_Weights.IMAGENET1K\_V1} from \texttt{torchvision} and is
loaded once per condition. Dataset bytes are the
\texttt{uoft-cs/cifar10} HuggingFace release at the revision pinned in
the released \texttt{requirements.txt}. A single command,
\texttt{python experiment.py}, on a vanilla RunPod \texttt{pytorch:2.4.0-cuda12.4.1}
image, reproduces the full nine-training sweep and rewrites
\texttt{summary.json} from scratch.


\section{Results}
\label{sec:results}

\subsection{Main result}

Across three random seeds, the predicted strict ordering
\textsc{Scratch} $<$ \textsc{Linear Probe} $<$ \textsc{Fine-tune} on
CIFAR-10 test top-1 accuracy holds at every seed and in the mean.
A 5-epoch ImageNet-pretrained linear probe attains
$\mathbf{0.8674 \pm 0.0016}$,
which is 8.44 percentage points
above an otherwise-identical ResNet-18 trained from random initialisation
($0.7830 \pm 0.0244$).
Fully fine-tuning the same ImageNet checkpoint end-to-end attains
$\mathbf{0.9119 \pm 0.0025}$,
a further 4.45 percentage points
on top of the linear probe and 12.90
percentage points above the from-scratch baseline. Equivalently,
fine-tuning reduces the from-scratch test error from
21.70\% to
8.81\%, a relative error reduction of
59.4\%. Figure~\ref{fig:main-bar}
visualises the three numbers side by side.

\begin{figure}[t]
  \centering
  \includegraphics[width=\linewidth]{figures/fig-main-bar.png}
  \caption{CIFAR-10 test top-1 accuracy by transfer-learning
    strategy, mean $\pm$ standard deviation across three random seeds.
    Scratch: $0.7830 \pm 0.0244$.
    Linear probe: $0.8674 \pm 0.0016$.
    Fine-tune: $0.9119 \pm 0.0025$.
    The strict ordering Scratch $<$ Linear $<$ Fine-tune is preserved
    on every individual seed.}
  \label{fig:main-bar}
\end{figure}

Table~\ref{tab:transfer-main} summarises the three conditions
together with their trainable-parameter counts. The pretrained
conditions exhibit very tight across-seed dispersion (linear probe
$\sigma=0.0016$, fine-tune
$\sigma=0.0025$), whereas the from-scratch
condition is approximately an order of magnitude noisier
($\sigma=0.0244$). This is consistent with
pretrained initialisations occupying a flatter region of the loss
landscape than random initialisation under the same 5-epoch budget.

\begin{table*}[!t]
\centering
\begin{tabular}{lcc}
\toprule
Condition & Trainable params & Test top-1 accuracy ($\uparrow$) \\
\midrule
ResNet-18 from scratch                       & $\sim$11.2\,M & $0.7830 \pm 0.0244$ \\
Linear probe (frozen backbone)               & 5{,}130        & $0.8674 \pm 0.0016$ \\
\textbf{Fine-tune (ImageNet ResNet-18)}      & $\sim$11.2\,M & $\mathbf{0.9119 \pm 0.0025}$ \\
\bottomrule
\end{tabular}
\caption{CIFAR-10 test top-1 accuracy on the official 10{,}000-image
test split, mean $\pm$ standard deviation across three random seeds
(0, 1, 2). All three conditions use the same ResNet-18 architecture
and 5-epoch training schedule on $224{\times}224$ inputs upsampled
from CIFAR-10's native $32{\times}32$. The linear probe trains only
the 10-class final layer on top of frozen ImageNet-pretrained
features; fine-tuning trains all parameters from the ImageNet
checkpoint.}
\label{tab:transfer-main}
\end{table*}

\subsection{Per-seed agreement and learning dynamics}

The strict ordering holds at the level of every individual seed, not
just on the mean: linear-probe accuracy exceeds the from-scratch
number on each of seeds 0, 1, and 2 (gaps of 9.17, 5.35, 10.81
percentage points respectively), and fine-tune exceeds linear probe
on each seed (gaps of 4.18, 4.38, 4.80
percentage points). No seed exhibits a non-monotone pattern.

The per-epoch test-accuracy trajectories
(Figure~\ref{fig:test-curves}) make the source of the gap explicit:
the linear probe converges within two epochs to its asymptotic
accuracy, fine-tuning continues to improve through epoch 5 and would
likely benefit from a longer schedule, while from-scratch is still
rising at epoch 5 and is the only condition for which the 5-epoch
budget appears to be the limiting factor. The training-loss curves
(Figure~\ref{fig:loss-curves}) show the same pattern from the
optimisation side: fine-tuning attains the lowest training loss, the
linear probe plateaus at a higher loss because the backbone cannot
adapt, and the from-scratch run is still descending when training
stops.

\begin{figure}[t]
  \centering
  \includegraphics[width=\linewidth]{figures/fig-test-curves.png}
  \caption{Per-epoch CIFAR-10 test top-1 accuracy for each
    transfer-learning strategy, mean $\pm$ standard deviation across
    three seeds. The linear probe converges within two epochs;
    fine-tuning continues improving through epoch 5; from-scratch is
    still rising at epoch 5.}
  \label{fig:test-curves}
\end{figure}

\begin{figure}[t]
  \centering
  \includegraphics[width=\linewidth]{figures/fig-loss-curves.png}
  \caption{Per-epoch training cross-entropy loss, mean $\pm$ standard
    deviation across three seeds. Fine-tuning attains the lowest
    training loss because it can adapt the backbone; the linear probe
    plateaus at a higher loss because the backbone is frozen;
    from-scratch is still descending when training stops.}
  \label{fig:loss-curves}
\end{figure}

\subsection{Compute footprint and limitations}

Each iteration trained nine ResNet-18 models (three conditions
$\times$ three seeds) on a single NVIDIA RTX A4000 (16\,GB) using
$224{\times}224$ inputs, batch size 128, and Adam (scratch and
fine-tune, lr $=10^{-3}$) or SGD with momentum 0.9 (linear probe,
lr $=10^{-2}$). Peak GPU memory was 3.13\,GB per
condition and total wall-clock per seed averaged 478\,s for the
three back-to-back trainings, comfortably within a single-GPU
laptop-scale compute envelope.

We deliberately did not search hyperparameters and we trained for
only five epochs per condition. Five epochs is sufficient for the
pretrained conditions to converge but leaves the from-scratch
baseline in the middle of its trajectory; absolute from-scratch
accuracy on CIFAR-10 with proper schedules and longer training is
reported in the literature to exceed 90\,\%, whereas our 5-epoch
number is 0.7830. The reported gaps should
therefore be read as a quantification of how quickly each strategy
reaches a competent classifier under a tight compute budget, the
regime where the question of "how useful is transfer learning?" is
most often asked in practice, rather than as a comparison of
asymptotic accuracies.


\section{Conclusion}
\label{sec:conclusion}

We decomposed the ImageNet-to-CIFAR-10 transfer-learning win for a
single ResNet-18 into the value of reusing pretrained features
(quantified by a linear probe over a frozen backbone) and the value
of adapting those features (quantified by full end-to-end
fine-tuning of the same checkpoint), holding architecture, training
schedule, batch size, input pipeline, and seed budget identical
across the three conditions. Across three random seeds and a 5-epoch
budget on a single 16\,GB GPU, the from-scratch baseline reached
$0.7830 \pm 0.0244$ test top-1 accuracy on the
official CIFAR-10 test split, the linear probe over frozen ImageNet
features reached $0.8674 \pm 0.0016$, and full fine-tuning
reached $\mathbf{0.9119 \pm 0.0025}$. The two
contributions split as $+8.44$ percentage points of feature
reuse and an additional $+4.45$ percentage points of
feature adaptation; the strict ordering held on every individual
seed.

Three limitations should temper any extrapolation. First, only one
backbone (ResNet-18) and one source domain (ImageNet-1K) were
evaluated, so the headline percentages should be read as a single
data point rather than a general law; deeper backbones or
out-of-distribution sources can shift both gaps in either direction.
Second, every condition was trained for only five epochs, which is
sufficient for the pretrained conditions to converge but leaves the
from-scratch baseline still descending; absolute from-scratch
accuracy on CIFAR-10 with proper schedules and longer training is
known to exceed 90\,\% in the literature, so the 0.7830 we
report should not be quoted as a state-of-the-art number. Third, the
$32{\times}32$-to-$224{\times}224$ upsampling that lets ImageNet
weights be applied at all is itself a confound when the from-scratch
condition is interpreted as a CIFAR-10 baseline rather than as a
control for the transfer comparison.

Two follow-ups would tighten the picture. The most direct is a
per-block freeze sweep: progressively unfreezing residual stages from
the top down (head only, then head plus stage 4, then plus stage 3,
and so on) would interpolate between the linear-probe and
fine-tuning end-points and locate which stages contribute the most
adaptation budget. The second is a budget sweep over training
epochs: the 8.44 and 4.45
percentage-point gaps reported here are 5-epoch numbers, and a
30-epoch or 100-epoch sweep would let the from-scratch baseline catch
up and the fine-tune number saturate, recasting the same
decomposition as a function of compute. Both follow-ups fit
comfortably inside the same single-GPU envelope and the same release
artefacts that accompany this paper.


% flushend handles balance automatically
\FloatBarrier
\bibliographystyle{unsrtnat}
\bibliography{references}

\end{document}
